\documentclass[11pt]{article}
\usepackage{fontspec,amsmath,amssymb,amsthm,hyperref}
\usepackage[margin=1in]{geometry}
\setlength{\emergencystretch}{3em}
\title{Construct the stationary coupling of one actual Gaussian LMC step}
\author{LangevinLean mathematical exposition}
\date{}
\begin{document}\maketitle
\noindent Source module: \texttt{AverageLMC/GaussianCanonical.lean}. Compile this proof document with XeLaTeX or LuaLaTeX. The accompanying HTML pairs the mathematical argument with the exact Lean statements and records the scope of the formalization.


The theorem starts from an arbitrary finite-cost initial coupling and the original smooth strongly convex potential. It constructs the independent Brownian extension, the actual integral solution, its stationary Gibbs law, and the exact Gaussian output law. The only local auxiliary input is the coordinate smoothing theorem proved separately.



Let \(\gamma\) be any finite-cost coupling of \(\nu\) and \(\pi\), with \(C_\gamma^2=\int\|x-y\|^2\,d\gamma\). The canonical independent Brownian extension realizes the actual update


\[\vartheta^+=\vartheta-h\nabla f(\vartheta)+\sqrt{2h}\,\xi,\qquad\xi\sim N(0,I_p),\quad\xi\perp\!\!\!\perp\vartheta.\]
\[(1+mh)W_2^2(\nu^+,\pi)\le(1-mh)C_\gamma^2+(M_{\rm av}+m)h^2p+h^2(hL-1)\int\|\nabla f(x)\|^2\,d\nu(x).\]
\section{Construct, identify and estimate the actual transition}


Extend the initial coupling by the explicit canonical product Brownian probability space. Its initial algorithm state X has law \ensuremath{\mu} and its reference state Z has the true Gibbs law. The Lipschitz drift constructs a continuous L\ensuremath{{}^{2}} Langevin integral solution; stationarity follows from the already proved Gibbs-stationarity theorem. Independent scalar Brownian paths provide independent standard Gaussian innovations at every positive time. The algorithm center A=X−h∇f(X) is a measurable function of the initial state, so it remains independent of those innovations. The endpoint A+B(h) has exactly the ordinary Gaussian LMC transition law. Local coordinate smoothing at A gives the affine potential bound needed by the Gaussian energy theorem. Integrated strong convexity controls ‖A−Z‖\ensuremath{{}^{2}}, and integrated deterministic descent controls E f(A)−E f(X). These potential terms cancel. Identifying the lifted gradient norm with the score moment under \ensuremath{\mu} leaves precisely (1−mh) times the initial coupling cost squared, (Mav+m)h\ensuremath{{}^{2}}d, and h\ensuremath{{}^{2}}(hL−1) times the score second moment. Output transport finiteness follows from the constructed L\ensuremath{{}^{2}} coupling before any real-valued distance is used.



\noindent Lean declarations: \texttt{lmc\_coupling\_step}.

\end{document}
